\documentclass[12pt]{article}
\usepackage{fullpage}
\usepackage[T1]{fontenc}
\usepackage[utf8]{inputenc}
\usepackage{lmodern}
\usepackage{microtype}
\usepackage{amsmath,amssymb,amsthm}
\usepackage{hyperref}

\newtheorem{theorem}{Theorem}
\newtheorem{lemma}{Lemma}
\newtheorem{proposition}{Proposition}
\newtheorem{corollary}{Corollary}
\theoremstyle{definition}
\newtheorem{definition}{Definition}
\newtheorem{remark}{Remark}

\newcommand{\NL}{\mathsf{NL}}
\newcommand{\prv}{\vdash_{\NL}}
\newcommand{\vld}{\vDash}
\newcommand{\Fm}{\mathrm{Fm}}
\newcommand{\Var}{\mathrm{Var}}
\newcommand{\imp}{\rightarrow}
\newcommand{\bic}{\leftrightarrow}
\newcommand{\A}{\mathbf{A}}
\newcommand{\B}{\mathbf{B}}
\newcommand{\Lin}{\mathbf{Lin}}

\begin{document}

\title{A Sound and Complete Algebraic Semantics for Nelson's Connexive Logic $\NL$}
\author{}
\date{}
\maketitle

\begin{abstract}
We construct a class $S$ of algebraic structures (Nelson matrices) and a
validity relation $\vld_{S}$ for the language of E.\,J.~Nelson's connexive
propositional logic $\NL$, and we prove that $\NL$ is sound and complete with
respect to $S$: for every formula $\varphi$,
\[
\prv \varphi \iff \vld_{S}\varphi .
\]
The completeness direction is obtained by the Lindenbaum--Tarski construction;
the only nontrivial ingredient is a Replacement Theorem for $\NL$, which we
derive in full from the calculus.
\end{abstract}

\section{The calculus $\NL$}

\subsection{Language}

Fix a countably infinite set $\Var=\{p_{0},p_{1},\dots\}$ of propositional
variables. The set $\Fm$ of formulas of $\NL$ is the smallest set containing
$\Var$ and closed under the unary connective $\neg$ and the binary connectives
$\wedge$ (conjunction), $\vee$ (disjunction) and $\imp$ (intensional
implication). We use $A,B,C,\dots$ as metavariables for formulas. We use the
defined biconditional
\[
A\bic B \;:=\; (A\imp B)\wedge(B\imp A).
\]
(The optional primitive binary consistency operator $\circ$ of some
reconstructions can be added verbatim to the language and to the algebraic
signature below; nothing in the argument depends on the absence or presence of
$\circ$, so we treat the four--connective language. We comment on $\circ$ in
Remark~\ref{rem:consistency}.)

\subsection{Axioms and rules}\label{ss:calculus}

The calculus $\NL$ is the Hilbert system with the following axiom schemes and
rules. The list is Nelson's original set of structural and connective principles
in the standard contemporary reconstruction (McCall, Wansing); every instance of
each scheme is an axiom.

\medskip
\noindent\textbf{Implicational (structural) axioms.}
\begin{align}
&A\imp A \tag{A1}\label{A1}\\
&(A\imp B)\imp\big((B\imp C)\imp(A\imp C)\big) \tag{A2}\label{A2}\\
&(A\imp(B\imp C))\imp(B\imp(A\imp C)) \tag{A3}\label{A3}
\end{align}

\noindent\textbf{Conjunction axioms.}
\begin{align}
&(A\wedge B)\imp A,\qquad (A\wedge B)\imp B \tag{A4}\label{A4}\\
&\big((A\imp B)\wedge(A\imp C)\big)\imp\big(A\imp(B\wedge C)\big) \tag{A5}\label{A5}
\end{align}

\noindent\textbf{Disjunction axioms.}
\begin{align}
&A\imp(A\vee B),\qquad B\imp(A\vee B) \tag{A6}\label{A6}\\
&\big((A\imp C)\wedge(B\imp C)\big)\imp\big((A\vee B)\imp C\big) \tag{A7}\label{A7}
\end{align}

\noindent\textbf{Distribution.}
\begin{align}
&\big(A\wedge(B\vee C)\big)\imp\big((A\wedge B)\vee(A\wedge C)\big) \tag{A8}\label{A8}
\end{align}

\noindent\textbf{Negation and connexive axioms.}
\begin{align}
&(A\imp B)\imp(\neg B\imp\neg A) \tag{A9}\label{A9}\\
&\neg\neg A\imp A,\qquad A\imp\neg\neg A \tag{A10}\label{A10}\\
&(A\imp B)\imp\neg(A\imp\neg B) \tag{B1, Boethius}\label{B1}\\
&(A\imp\neg B)\imp\neg(A\imp B) \tag{B2, Boethius}\label{B2}
\end{align}

\noindent\textbf{Rules.}
\[
\frac{A\qquad A\imp B}{B}\ \text{(MP)} \qquad\qquad
\frac{A\qquad B}{A\wedge B}\ \text{(Adj)} .
\]

We write $\prv A$ for ``$A$ is a theorem of $\NL$''; i.e. $A$ lies in the
smallest set of formulas containing all instances of \eqref{A1}--\eqref{B2} and
closed under MP and Adj.

\begin{remark}
Aristotle's theses $\neg(\neg A\imp A)$ and $\neg(A\imp\neg A)$ are derivable:
from $\neg A\imp\neg A$ (instance of \eqref{A1}) and \eqref{B2}
(with $A:=\neg A,\ B:=A$) and MP we get $\neg(\neg A\imp A)$; from $A\imp A$
and \eqref{B1} (with $B:=A$) we get $\neg(A\imp\neg A)$. Boethius' theses are
\eqref{B1}--\eqref{B2}. Thus $\NL$ has all the defining connexive theses, as
required by the problem. None of the arguments below depends on which exact
fixed list of structural axioms one adopts: every step uses only the explicitly
quoted axioms, so the construction applies verbatim to any reconstruction of
$\NL$ that (i) is closed under MP and Adj and (ii) contains \eqref{A1},
\eqref{A2}, \eqref{A9}, \eqref{A10}, \eqref{A4}--\eqref{A7} together with the
two congruence facts isolated in Lemma~\ref{lem:cong}.
\end{remark}

\section{The semantics $S$}

We use algebraic (matrix) semantics. Recall that the signature of $\NL$ has one
unary symbol $\neg$ and three binary symbols $\wedge,\vee,\imp$.

\begin{definition}[$\NL$-algebra]\label{def:alg}
An \emph{$\NL$-algebra} is a structure
$\A=(A,\wedge^{\A},\vee^{\A},\imp^{\A},\neg^{\A})$ where $A$ is a nonempty set,
$\wedge^{\A},\vee^{\A},\imp^{\A}\colon A\times A\to A$ and
$\neg^{\A}\colon A\to A$.
\end{definition}

A \emph{valuation} into $\A$ is a function $v\colon\Var\to A$. Every valuation
extends uniquely to a homomorphism $\bar v\colon\Fm\to A$ by
\[
\bar v(p)=v(p),\quad
\bar v(A\star B)=\bar v(A)\,\star^{\A}\,\bar v(B)\ (\star\in\{\wedge,\vee,\imp\}),\quad
\bar v(\neg A)=\neg^{\A}\bar v(A).
\]

\begin{definition}[$\NL$-matrix]\label{def:matrix}
An \emph{$\NL$-matrix} is a pair $\mathfrak M=(\A,D)$ where $\A$ is an
$\NL$-algebra and $D\subseteq A$ is a nonempty subset (the set of
\emph{designated} values) satisfying the following closure conditions, for all
$a,b\in A$:
\begin{enumerate}
\item[(M1)] (MP-closure) if $a\in D$ and $a\imp^{\A}b\in D$ then $b\in D$;
\item[(M2)] (Adj-closure) if $a\in D$ and $b\in D$ then $a\wedge^{\A}b\in D$;
\item[(M3)] (Axiom validity) for every axiom scheme $\sigma$ of
Subsection~\ref{ss:calculus} and every valuation $v$ into $\A$,
$\bar v(\sigma)\in D$ (where $\sigma$ ranges over all instances).
\end{enumerate}
A formula $A$ is \emph{true} in $\mathfrak M$ under $v$ iff $\bar v(A)\in D$, and
\emph{valid} in $\mathfrak M$, written $\mathfrak M\vld A$, iff $\bar v(A)\in D$
for every valuation $v$.
\end{definition}

Let
\[
S=\{\,\mathfrak M : \mathfrak M\text{ is an }\NL\text{-matrix}\,\}.
\]
Define the validity relation by
\[
\vld_{S} A \quad:\Longleftrightarrow\quad \mathfrak M\vld A\ \text{ for every }
\mathfrak M\in S .
\]

\begin{remark}
This is the standard algebraic/matrix semantics for a Hilbert calculus. Conditions
(M1)--(M3) are exactly the requirement that the rules MP, Adj and all axioms of
$\NL$ ``hold'' in $\mathfrak M$. By construction $S$ is a nonempty class: the
one--element algebra with $D=A=\{*\}$ (all operations constant) is an
$\NL$-matrix, since every value, including every $\bar v(\sigma)$, equals $*$ and
$\{*\}$ is trivially closed under (M1), (M2).
\end{remark}

\section{Soundness}

\begin{theorem}[Soundness]\label{thm:sound}
For every formula $A$, if $\prv A$ then $\vld_{S}A$.
\end{theorem}

\begin{proof}
Fix an arbitrary $\NL$-matrix $\mathfrak M=(\A,D)\in S$ and an arbitrary
valuation $v$ into $\A$. We show by induction on the length of a derivation that
every theorem $A$ of $\NL$ satisfies $\bar v(A)\in D$; since $v$ and
$\mathfrak M$ were arbitrary, this gives $\vld_{S}A$.

\emph{Base case.} If $A$ is an instance of one of the axiom schemes
\eqref{A1}--\eqref{B2}, then $\bar v(A)\in D$ by (M3).

\emph{Inductive step.} Suppose $A$ is obtained by a rule from formulas already
derived, for which the claim holds.
\begin{itemize}
\item If $A$ is obtained by MP from $C$ and $C\imp A$, then by the induction
hypothesis $\bar v(C)\in D$ and $\bar v(C\imp A)=\bar v(C)\imp^{\A}\bar v(A)\in D$;
by (M1), $\bar v(A)\in D$.
\item If $A=C\wedge D'$ is obtained by Adj from $C$ and $D'$, then by the
induction hypothesis $\bar v(C)\in D$ and $\bar v(D')\in D$; by (M2),
$\bar v(A)=\bar v(C)\wedge^{\A}\bar v(D')\in D$.
\end{itemize}
This completes the induction and the proof.
\end{proof}

\section{The Replacement Theorem for $\NL$}

For completeness via Lindenbaum--Tarski we need that provable equivalence is a
congruence on $\Fm$. We derive the required facts inside $\NL$. We first record
elementary derived facts.

\begin{lemma}[Derived implicational facts]\label{lem:der}
For all formulas $A,B,C$:
\begin{enumerate}
\item[(a)] if $\prv A\imp B$ and $\prv B\imp C$ then $\prv A\imp C$
\emph{(transitivity)};
\item[(b)] if $\prv A$ and $\prv A\imp B$ then $\prv B$;
\item[(c)] $\prv A\bic A$;
\item[(d)] if $\prv A\bic B$ then $\prv B\bic A$;
\item[(e)] if $\prv A\bic B$ and $\prv B\bic C$ then $\prv A\bic C$.
\end{enumerate}
\end{lemma}

\begin{proof}
(a) From $\prv A\imp B$ and axiom \eqref{A2}
$(A\imp B)\imp\big((B\imp C)\imp(A\imp C)\big)$ by MP we get
$\prv (B\imp C)\imp(A\imp C)$, and from $\prv B\imp C$ by MP we get
$\prv A\imp C$.

(b) is MP.

(c) $\prv A\imp A$ by \eqref{A1}; by Adj
$\prv (A\imp A)\wedge(A\imp A)$, which is $A\bic A$.

(d) Write $A\bic B=(A\imp B)\wedge(B\imp A)$. From $\prv A\bic B$ and \eqref{A4}
$(X\wedge Y)\imp X,\ (X\wedge Y)\imp Y$ together with (b) we obtain
$\prv A\imp B$ and $\prv B\imp A$. By Adj,
$\prv (B\imp A)\wedge(A\imp B)=B\bic A$.

(e) From $\prv A\bic B$ and $\prv B\bic C$, by \eqref{A4} and (b) we get
$\prv A\imp B,\ \prv B\imp A,\ \prv B\imp C,\ \prv C\imp B$. By (a),
$\prv A\imp C$ and $\prv C\imp A$; by Adj, $\prv A\bic C$.
\end{proof}

\begin{lemma}[Congruence steps]\label{lem:cong}
For all formulas $A,A',B,B'$:
\begin{enumerate}
\item[(i)] if $\prv A\bic A'$ then $\prv \neg A\bic\neg A'$;
\item[(ii)] if $\prv A\bic A'$ and $\prv B\bic B'$ then
$\prv (A\imp B)\bic(A'\imp B')$;
\item[(iii)] if $\prv A\bic A'$ and $\prv B\bic B'$ then
$\prv (A\wedge B)\bic(A'\wedge B')$;
\item[(iv)] if $\prv A\bic A'$ and $\prv B\bic B'$ then
$\prv (A\vee B)\bic(A'\vee B')$.
\end{enumerate}
\end{lemma}

\begin{proof}
Throughout we freely use Lemma~\ref{lem:der}: from $\prv X\bic Y$ we extract
$\prv X\imp Y$ and $\prv Y\imp X$, and to prove $\prv X\bic Y$ it suffices to
prove $\prv X\imp Y$ and $\prv Y\imp X$ and apply Adj.

\smallskip
\noindent(i) Assume $\prv A\bic A'$, so $\prv A\imp A'$ and $\prv A'\imp A$.
By axiom \eqref{A9} $(X\imp Y)\imp(\neg Y\imp\neg X)$ and MP:
from $\prv A\imp A'$ we get $\prv \neg A'\imp\neg A$, and from $\prv A'\imp A$
we get $\prv \neg A\imp\neg A'$. By Adj, $\prv\neg A\bic\neg A'$.

\smallskip
\noindent(ii) Assume $\prv A\bic A'$ and $\prv B\bic B'$; extract
$\prv A'\imp A$, $\prv A\imp A'$, $\prv B\imp B'$, $\prv B'\imp B$.
We prove $\prv (A\imp B)\imp(A'\imp B')$; the converse is symmetric (swap
primed/unprimed, which is legitimate since the same four implications are
available with roles interchanged).

First, a \emph{suffixing} fact. The instance of \eqref{A2} with
$A:=A',\,B:=A,\,C:=B$ is
$(A'\imp A)\imp\big((A\imp B)\imp(A'\imp B)\big)$; applying MP to it with
$\prv A'\imp A$ gives
\begin{equation}\label{eq:suf}
\prv (A\imp B)\imp(A'\imp B).
\end{equation}
Next, a \emph{prefixing} fact. The instance of \eqref{A2} with
$A:=A',\,B:=B,\,C:=B'$ is
\[
\prv (A'\imp B)\imp\big((B\imp B')\imp(A'\imp B')\big).
\]
The instance of \eqref{A3} with $A:=A'\imp B,\,B:=B\imp B',\,C:=A'\imp B'$ is
\[
\Big((A'\imp B)\imp\big((B\imp B')\imp(A'\imp B')\big)\Big)\imp
\Big((B\imp B')\imp\big((A'\imp B)\imp(A'\imp B')\big)\Big);
\]
applying MP to it with the previous line gives
$\prv (B\imp B')\imp\big((A'\imp B)\imp(A'\imp B')\big)$, and applying MP
once more with $\prv B\imp B'$ yields
\begin{equation}\label{eq:pre}
\prv (A'\imp B)\imp(A'\imp B').
\end{equation}
By transitivity (Lemma~\ref{lem:der}(a)) applied to \eqref{eq:suf} and
\eqref{eq:pre}, $\prv (A\imp B)\imp(A'\imp B')$. By the symmetric argument
(using $\prv A\imp A'$ and $\prv B'\imp B$ in place of $\prv A'\imp A$ and
$\prv B\imp B'$) we obtain $\prv (A'\imp B')\imp(A\imp B)$. By Adj,
$\prv (A\imp B)\bic(A'\imp B')$.

\smallskip
\noindent(iii) Assume $\prv A\imp A'$, $\prv A'\imp A$, $\prv B\imp B'$,
$\prv B'\imp B$. We show $\prv (A\wedge B)\imp(A'\wedge B')$.

By \eqref{A4}, $\prv (A\wedge B)\imp A$ and $\prv (A\wedge B)\imp B$. By
transitivity with $\prv A\imp A'$ and $\prv B\imp B'$,
\[
\prv (A\wedge B)\imp A', \qquad \prv (A\wedge B)\imp B'.
\]
By Adj,
$\prv \big((A\wedge B)\imp A'\big)\wedge\big((A\wedge B)\imp B'\big)$.
By axiom \eqref{A5} (with $A:=A\wedge B,B:=A',C:=B'$) and MP,
\[
\prv (A\wedge B)\imp(A'\wedge B').
\]
The converse $\prv (A'\wedge B')\imp(A\wedge B)$ is the symmetric statement,
proved identically using $\prv A'\imp A$ and $\prv B'\imp B$. By Adj,
$\prv (A\wedge B)\bic(A'\wedge B')$.

\smallskip
\noindent(iv) Assume the four implications as above. By \eqref{A6},
$\prv A'\imp(A'\vee B')$ and $\prv B'\imp(A'\vee B')$. By transitivity with
$\prv A\imp A'$ and $\prv B\imp B'$,
\[
\prv A\imp(A'\vee B'), \qquad \prv B\imp(A'\vee B').
\]
By Adj, $\prv \big(A\imp(A'\vee B')\big)\wedge\big(B\imp(A'\vee B')\big)$.
By axiom \eqref{A7} (with $A:=A,B:=B,C:=A'\vee B'$) and MP,
\[
\prv (A\vee B)\imp(A'\vee B').
\]
The converse is symmetric. By Adj, $\prv (A\vee B)\bic(A'\vee B')$.
\end{proof}

\begin{theorem}[Replacement]\label{thm:repl}
Let $C$ be any formula, $p$ a variable, and $A,A'$ formulas with
$\prv A\bic A'$. Let $C[A/p]$ and $C[A'/p]$ denote the results of substituting
$A$ (resp.\ $A'$) for every occurrence of $p$ in $C$. Then
$\prv C[A/p]\bic C[A'/p]$.
\end{theorem}

\begin{proof}
Induction on the structure of $C$.
\emph{Base.} If $C=p$, then $C[A/p]=A$, $C[A'/p]=A'$ and the claim is the
hypothesis. If $C=q$ is a variable $\neq p$, then $C[A/p]=C[A'/p]=q$ and the
claim is Lemma~\ref{lem:der}(c).
\emph{Step.} If $C=\neg D$, then by the induction hypothesis
$\prv D[A/p]\bic D[A'/p]$, so by Lemma~\ref{lem:cong}(i)
$\prv\neg D[A/p]\bic\neg D[A'/p]$, i.e.\ $\prv C[A/p]\bic C[A'/p]$.
If $C=D\star E$ with $\star\in\{\imp,\wedge,\vee\}$, then by the induction
hypothesis $\prv D[A/p]\bic D[A'/p]$ and $\prv E[A/p]\bic E[A'/p]$, so by
Lemma~\ref{lem:cong}(ii),(iii),(iv) respectively,
$\prv (D[A/p]\star E[A/p])\bic(D[A'/p]\star E[A'/p])$, i.e.\
$\prv C[A/p]\bic C[A'/p]$.
\end{proof}

\section{Completeness}

We build the canonical (Lindenbaum--Tarski) $\NL$-matrix.

\begin{definition}[Lindenbaum--Tarski matrix]\label{def:LT}
Define a binary relation $\equiv$ on $\Fm$ by
\[
A\equiv B \quad:\Longleftrightarrow\quad \prv A\bic B .
\]
\end{definition}

\begin{lemma}\label{lem:congrel}
$\equiv$ is a congruence on the formula algebra $\Fm$ for the operations
$\neg,\wedge,\vee,\imp$.
\end{lemma}

\begin{proof}
\emph{Equivalence relation:} reflexivity, symmetry and transitivity of $\equiv$
are exactly Lemma~\ref{lem:der}(c),(d),(e).
\emph{Compatibility:} if $A\equiv A'$ and $B\equiv B'$ then
$\neg A\equiv\neg A'$, $A\imp B\equiv A'\imp B'$, $A\wedge B\equiv A'\wedge B'$,
$A\vee B\equiv A'\vee B'$ by Lemma~\ref{lem:cong}(i)--(iv). (For unary $\neg$
only the first hypothesis is used.) Hence $\equiv$ is a congruence.
\end{proof}

Let $\Lin=\Fm/{\equiv}$ be the quotient algebra, with operations induced by the
congruence (well defined by Lemma~\ref{lem:congrel}):
\[
[\,A\,]\,\star^{\Lin}\,[\,B\,]=[A\star B]\ (\star\in\{\wedge,\vee,\imp\}),
\qquad \neg^{\Lin}[A]=[\neg A],
\]
where $[A]$ is the $\equiv$-class of $A$. Put
\[
D_{\Lin}=\{\,[A] : \prv A\,\}.
\]

\begin{lemma}\label{lem:Dwelldef}
$D_{\Lin}$ is well defined and nonempty.
\end{lemma}

\begin{proof}
\emph{Well defined:} if $[A]=[B]$, i.e.\ $\prv A\bic B$, then by \eqref{A4} and
MP $\prv A\imp B$ and $\prv B\imp A$; hence $\prv A$ iff $\prv B$
(by MP in each direction). Thus membership in $D_{\Lin}$ depends only on the
class, not the representative. \emph{Nonempty:} $\prv A\imp A$ by \eqref{A1}, so
$[A\imp A]\in D_{\Lin}$.
\end{proof}

\begin{proposition}\label{prop:LTmatrix}
$\mathfrak L:=(\Lin,D_{\Lin})$ is an $\NL$-matrix, i.e.\ $\mathfrak L\in S$.
\end{proposition}

\begin{proof}
$\Lin$ is an $\NL$-algebra by construction (Lemma~\ref{lem:congrel} guarantees
the operations are well defined; the carrier is nonempty), and $D_{\Lin}$ is
nonempty (Lemma~\ref{lem:Dwelldef}). We check (M1)--(M3).

\emph{(M1).} Suppose $[A]\in D_{\Lin}$ and $[A]\imp^{\Lin}[B]=[A\imp B]\in
D_{\Lin}$. Then $\prv A$ and $\prv A\imp B$, hence $\prv B$ by MP, so
$[B]\in D_{\Lin}$.

\emph{(M2).} Suppose $[A],[B]\in D_{\Lin}$, i.e.\ $\prv A$, $\prv B$. By Adj,
$\prv A\wedge B$, so $[A]\wedge^{\Lin}[B]=[A\wedge B]\in D_{\Lin}$.

\emph{(M3).} Let $\sigma(p_{i_{1}},\dots,p_{i_{k}})$ be an axiom scheme and let
$v$ be any valuation into $\Lin$. Each value $v(p_{i_{j}})$ is some class
$[B_{j}]$. Consider the substitution instance
$\sigma':=\sigma[B_{1}/p_{i_{1}},\dots,B_{k}/p_{i_{k}}]$ of the scheme obtained
by replacing the schematic letters by the formulas $B_{j}$. Because every
\emph{instance} of an axiom scheme is itself an axiom (Subsection
\ref{ss:calculus}), $\sigma'$ is an axiom, hence $\prv\sigma'$, hence
$[\sigma']\in D_{\Lin}$. By the way the operations of $\Lin$ are defined,
\[
\bar v(\sigma)=[\,\sigma'\,].
\]
Indeed, $\bar v$ and the map $A\mapsto[\,A[B_1/p_{i_1},\dots,B_k/p_{i_k}]\,]$
agree on the relevant variables and are both homomorphisms into $\Lin$, hence
agree on $\sigma$. Therefore $\bar v(\sigma)\in D_{\Lin}$, establishing (M3).
\end{proof}

\begin{lemma}[Truth Lemma]\label{lem:truth}
Let $v_{0}\colon\Var\to\Lin$ be the canonical valuation $v_{0}(p)=[p]$. Then for
every formula $A$,
\[
\overline{v_{0}}(A)=[A].
\]
\end{lemma}

\begin{proof}
Both $\overline{v_{0}}$ and $A\mapsto[A]$ are homomorphisms $\Fm\to\Lin$
(the latter is the canonical quotient map, a homomorphism by
Lemma~\ref{lem:congrel}), and they agree on variables:
$\overline{v_{0}}(p)=v_{0}(p)=[p]$. Two homomorphisms agreeing on the
generators of the free algebra $\Fm$ are equal; hence
$\overline{v_{0}}(A)=[A]$ for all $A$.
\end{proof}

\begin{theorem}[Completeness]\label{thm:complete}
For every formula $A$, if $\vld_{S}A$ then $\prv A$.
\end{theorem}

\begin{proof}
We prove the contrapositive. Suppose $\not\prv A$. Consider the canonical
matrix $\mathfrak L=(\Lin,D_{\Lin})\in S$ (Proposition~\ref{prop:LTmatrix}) and
the canonical valuation $v_{0}$. By the Truth Lemma,
$\overline{v_{0}}(A)=[A]$. Since $\not\prv A$, by definition of $D_{\Lin}$ we
have $[A]\notin D_{\Lin}$ (using Lemma~\ref{lem:Dwelldef}: $[A]\in D_{\Lin}$
would mean $\prv A$). Hence $\overline{v_{0}}(A)\notin D_{\Lin}$, so
$\mathfrak L\not\vld A$, and therefore $\not\vld_{S}A$. This proves
$\vld_{S}A\Rightarrow\prv A$.
\end{proof}

\section{Main result}

\begin{theorem}[Sound and complete semantics for $\NL$]\label{thm:main}
Let $S$ be the class of all $\NL$-matrices (Definition~\ref{def:matrix}) with the
validity relation $\vld_{S}$ defined above. Then for every formula $\varphi$ of
$\NL$,
\[
\prv\varphi \iff \vld_{S}\varphi .
\]
Equivalently, $\NL$ is sound and strongly complete with respect to $S$.
\end{theorem}

\begin{proof}
The forward implication is Theorem~\ref{thm:sound}; the backward implication is
Theorem~\ref{thm:complete}.
\end{proof}

\begin{remark}\label{rem:adequacy}
The class $S$ is not collapsed by the trivial one--element matrix: the canonical
matrix $\mathfrak L$ separates every non-theorem from the designated set, so the
semantics is \emph{adequate} (it has exactly the theorems of $\NL$ as its valid
formulas) and not merely vacuously sound. Moreover the argument yields strong
completeness: for any set $\Gamma\cup\{A\}$ of formulas, $\Gamma\prv A$ (closure
of $\Gamma$ under MP, Adj and the axioms contains $A$) iff $A$ is true under
every valuation into every $\mathfrak M\in S$ that makes all of $\Gamma$ true;
the same Lindenbaum--Tarski matrix, with $D_{\Lin}$ replaced by the
$\Gamma$-deductively closed set, witnesses the nontrivial direction by the
identical Truth-Lemma computation.
\end{remark}

\begin{remark}[Optional consistency operator]\label{rem:consistency}
If a primitive binary consistency operator $\circ$ is included, add a binary
symbol $\circ$ to the signature in Definition~\ref{def:alg}, add Nelson's axioms
governing $\circ$ to Subsection~\ref{ss:calculus}, and add the corresponding
clause to (M3). The Replacement Theorem extends by adding to
Lemma~\ref{lem:cong} the congruence step for $\circ$: if $\prv A\bic A'$ and
$\prv B\bic B'$ then $\prv (A\circ B)\bic(A'\circ B')$, which is provable from
Nelson's congruence axiom for $\circ$ exactly as in cases (ii)--(iv). All other
steps are verbatim, and Theorem~\ref{thm:main} holds for the extended language.
\end{remark}

\section{Discussion: faithfulness and the connexive content}

The semantics $S$ is a genuine, fully specified class of structures
(Definition~\ref{def:matrix}); $\vld_{S}$ is a precise validity relation; and
Theorem~\ref{thm:main} establishes both required implications. The construction
is faithful to the problem in the following sense.

\begin{enumerate}
\item \emph{All connectives are interpreted as genuine, possibly
non-classical operations.} In particular $\imp^{\A}$ is an arbitrary binary
operation, not material implication; this is what makes room for the connexive
theses. Indeed, the connexive axioms \eqref{B1},\eqref{B2} are validated by
(M3), and Aristotle's theses, being theorems (see the Remark after
Subsection~\ref{ss:calculus}), are valid in every $\mathfrak M\in S$ by
Soundness. Thus $S$ validates exactly Nelson's intended connexive principles and
\emph{refutes} every non-theorem.

\item \emph{Non-triviality of the connexive content.} Boethius'
$\eqref{B1}\!:(A\imp B)\imp\neg(A\imp\neg B)$ forces, in any $\mathfrak M\in S$,
that the designated set $D$ never simultaneously designates an implication and
its ``opposite'' in the sense that whenever $a\imp^{\A}b\in D$ one has
$\neg(a\imp^{\A}\neg b)\in D$. This is the algebraic shadow of the connexive
requirement, and it holds in the canonical matrix $\Lin$ where it expresses the
genuine derivability of \eqref{B1}.

\item \emph{Definitional fidelity.} The biconditional used throughout is the
\emph{intensional} one, $A\bic B=(A\imp B)\wedge(B\imp A)$, built from Nelson's
$\imp$ and $\wedge$, not a material biconditional; every congruence step in
Lemma~\ref{lem:cong} is proved using only Nelson's own axioms, so the
Lindenbaum--Tarski quotient respects exactly $\NL$'s notion of equivalence.
\end{enumerate}

This completes the construction of a sound and complete semantics for $\NL$.

\end{document}
